% Generated by eval/scripts/write_collinearity_explanation.py.
The single-predictor fits ask how the estimate changes when the other substantive predictor is omitted. We fit distance without SPEC, and SPEC without distance, retaining model and attack-language indicators in both cases. These pooled comparisons use the same 32\,090 scored responses and the same standardization. The following comparison therefore reflects a change in adjustment, rather than a change in sample:

\begin{center}
\small
\begin{tabular}{lrr}
\toprule
Predictor & Distance and SPEC together & Without the other predictor \\
\midrule
Distance & $-0.112$ & $-0.385$ \\
SPEC & $0.954$ & $1.157$ \\
\bottomrule
\end{tabular}
\end{center}

Both coefficients keep their signs, but their sizes change: the distance estimate becomes more negative without SPEC, while the SPEC estimate becomes more positive without distance. The estimated size of an association therefore depends on what else the regression accounts for. These coefficients are on the logit of expected score; they are not changes of that size in the original score. Table~\ref{tab:graded-single} reports the single-predictor intervals and the IF/CONS sensitivities.

Omitting distance does not remove all overlap: SPEC's VIF is still 7.99 because it remains associated with the retained model/language indicators. Likewise, the decomposition has IF VIF 38.41 and CONS VIF 5.04. The simpler fits reveal sensitivity to adjustment; they do not isolate a causal effect or resolve omitted-variable bias.

For H\textsubscript{3}, the implication is that ranking distance and SPEC as independent explanations requires caution. Standardization makes their coefficient units comparable, but the estimates still depend on overlapping predictors and cohort choices. We therefore consider their signs, sizes, intervals, and stability together. A larger absolute coefficient or smaller $p$-value alone does not establish a stronger and more stable predictor.

Finally, convergence answers a different question: whether the fitting algorithm reached a solution. All reported fits pass convergence, finite-estimate, and positive-definite covariance checks. Those checks make the estimates numerically usable; a successfully fitted model can still have substantial collinearity.
